\clearpage
\section{Rural--Urban Differences and Development}\label{sec:additional_rural_urban}
\suppressfloats[t]

The rural point estimate is lower than the urban estimate, but the difference is imprecise and sensitive to which cities are included. Figure~\ref{fig:additional_persistence} also reports the descriptive urban UP comparison; the meta-analysis here retains Mumbai, Delhi, Chandigarh, and Bengaluru as its urban settings. Comparisons within rural Rajasthan and UP also do not establish a consistent relationship between development and subsequent electoral gains.

\subsection{Rural--urban meta-analysis}

We test whether the prior-reservation contrast is smaller in rural settings using a meta-regression of setting-average persistence estimates on an intercept and a rural indicator. Mumbai, Delhi, Chandigarh, and Bengaluru are urban; Rajasthan, Uttar Pradesh, and West Bengal are rural. The model jointly estimates the subgroup means and their difference, using the same within-setting averages and assumed sampling variances as Section~\ref{sec:additional_pool}. It estimates a common residual between-setting variance by REML and uses modified Knapp--Hartung inference with the number of settings minus two degrees of freedom. The coefficient tests equality of the persistence contrasts, rather than equality of the underlying female-officeholding rates or separate significance within each subgroup.

The fitted rural and urban persistence contrasts are \appRuralMean{} and \appUrbanMean{} percentage points, respectively. Their difference, rural minus urban, is \appRuralGap{} points, with a 95\% interval of \appRuralGapInterval{} and a two-sided $p$-value of \appRuralGapP{}. The rural point estimate is lower, but the interval permits both a substantial rural disadvantage and an advantage. The test establishes neither a difference nor equivalence.

\subsection{Sensitivity of the rural--urban difference}\label{sec:rural_urban_pooling_sensitivity}

Table~\ref{tab:additional_pooling_methods} compares election weights and dependence assumptions, retaining all thirteen contrasts in seven settings. Except for the reference row, it uses the estimated cross-election covariance in Rajasthan and UP and varies the assumed correlation only in West Bengal and Chandigarh. Equal weights retain the average across observed elections. Generalized least-squares (GLS) weights instead minimize sampling variance under a common within-setting effect: for covariance matrix $V$, the weights are $a=V^{-1}\mathbf{1}/(\mathbf{1}'V^{-1}\mathbf{1})$ and the variance is $a'Va$. When effects differ across elections, GLS changes the target to a precision-weighted contrast. Both approaches retain the same random-effects pooling and modified Knapp--Hartung inference across geographic settings.

Changing the weights and the remaining correlation assumptions has little effect on either the pooled estimate or the rural--urban difference. The rural--urban intervals remain wide and include zero throughout.

\begin{table}[!htbp]
\centering
\caption{Alternative pooling methods for persistence and the rural--urban difference}
\label{tab:additional_pooling_methods}
\footnotesize
\setlength{\tabcolsep}{3pt}
\begin{tabular}{lcrrr}
\toprule
Election weights & Assumed correlation & Pooled [95\% interval] & Rural $-$ urban [95\% interval] & $p$ \\
\midrule
Equal (reference) & 1 everywhere & 1.5 [-1.5, 4.5] & -1.3 [-11.3, 8.7] & 0.745 \\
Equal & 1 & 1.4 [-1.7, 4.4] & -1.5 [-11.5, 8.6] & 0.721 \\
Equal & 0.5 & 1.4 [-1.6, 4.3] & -1.5 [-11.4, 8.5] & 0.719 \\
Equal & 0 & 1.4 [-1.5, 4.2] & -1.5 [-11.3, 8.3] & 0.711 \\
GLS & 0.5 & 1.5 [-1.2, 4.1] & -1.1 [-10.8, 8.6] & 0.784 \\
GLS & 0 & 1.5 [-1.1, 4.0] & -1.2 [-10.7, 8.4] & 0.765 \\
\bottomrule
\end{tabular}

\begin{minipage}{0.96\textwidth}\small
\emph{Notes:} Estimates and 95\% intervals are in percentage points. All rows use the same thirteen election contrasts in seven geographic settings. The reference row assumes correlation one within every setting. Other rows use estimated block-cluster covariance in Rajasthan and UP; the assumed correlation applies only to West Bengal and Chandigarh. GLS denotes generalized least squares; equal and GLS election weights are defined above. Across settings, all rows use inverse sampling-plus-heterogeneity variance weights, REML, and modified Knapp--Hartung inference, with six degrees of freedom for the pooled estimate and five for the rural--urban difference. The $p$-value tests a zero rural--urban difference. Neither method treats repeated elections as independent settings.
\end{minipage}
\end{table}

Table~\ref{tab:additional_rural_urban} applies the dependence and sample-composition changes defined above to the rural--urban coefficient. Adding Bihar supplies a fourth rural setting under a longer reservation regime. Omitting West Bengal restricts both groups to directly elected offices. Omitting each city in turn tests its influence, leaving three urban settings. The urban variance correction adds one rather than one-half to each binary-outcome cell when estimating sampling variances, retaining the raw differences. The earliest- and latest-election specifications assess temporal composition without counting repeated elections as independent settings.

The confidence interval includes zero in every reported specification. Completing both unresolved urban outcomes in all four possible ways gives rural--urban point estimates from \appRuralGapMissingLow{} to \appRuralGapMissingHigh{} points; each interval also includes zero. The evidence is compatible with lower rural persistence, but the small number of settings leaves the contrast imprecise. Geography also coincides with differences in the office being contested, reservation rules, and election periods. Neither the primary subgroup analysis nor these sensitivity checks identify an effect of urbanization.

\begin{table}[!htbp]
\centering
\caption{Direct tests of the rural--urban persistence difference}
\label{tab:additional_rural_urban}
\footnotesize
\setlength{\tabcolsep}{3pt}
\begin{tabular}{lrrr}
\toprule
Specification & Rural $-$ urban & 95\% interval & $p$ \\
\midrule
Primary: maximum within-setting correlation & -1.3 & [-11.3, 8.7] & 0.745 \\
Within-setting correlation = 0 & -1.5 & [-11.3, 8.3] & 0.711 \\
Within-setting correlation = 0.5 & -1.4 & [-11.3, 8.5] & 0.732 \\
Minimum feasible within-setting correlation & -3.3 & [-9.9, 3.3] & 0.255 \\
Estimated covariance in Rajasthan and UP & -1.5 & [-11.5, 8.6] & 0.721 \\
Urban variance correction = 1 & -2.2 & [-12.4, 8.0] & 0.598 \\
Include Bihar's two-election regime & -0.9 & [-9.6, 7.8] & 0.812 \\
Omit Rajasthan & -1.6 & [-16.3, 13.0] & 0.773 \\
Omit Uttar Pradesh & -3.3 & [-18.4, 11.8] & 0.578 \\
Omit Mumbai & 1.3 & [-7.4, 10.1] & 0.695 \\
Omit Delhi & -0.3 & [-14.6, 14.0] & 0.957 \\
Omit Chandigarh & -1.2 & [-13.5, 11.1] & 0.798 \\
Omit Bengaluru & -6.8 & [-17.7, 4.0] & 0.156 \\
Omit West Bengal & -2.7 & [-17.5, 12.2] & 0.641 \\
Earliest comparison per setting & -0.4 & [-9.3, 8.5] & 0.910 \\
Latest comparison per setting & -0.8 & [-11.4, 9.7] & 0.848 \\
\bottomrule
\end{tabular}

\begin{minipage}{0.96\textwidth}\small
\emph{Notes:} Differences and 95\% intervals are in percentage points; negative estimates indicate lower persistence in rural settings. The primary model has three rural and four urban settings. Each row fits an intercept and rural indicator with a common residual heterogeneity variance, REML estimation, and modified Knapp--Hartung $t(k-2)$ inference. Reported $p$-values are two-sided tests of a zero rural--urban difference. The correlation and sample changes match Table~\ref{tab:additional_meta}; adding Bihar supplies a fourth rural setting. Omitting each city is an influence diagnostic with three urban settings remaining. These exploratory comparisons do not identify a causal effect of urbanization.
\end{minipage}
\end{table}
\FloatBarrier

\subsection{Development within rural settings}\label{sec:app_het}

We interact prior quota assignment with three 2001 Census characteristics measured through SHRUG: proximity to towns, infrastructure availability, and female literacy. All three interactions enter the same model with outcome-year block fixed effects and clustering.

Near Town indicates a GP at or below the median distance to the nearest town. High Infrastructure indicates an above-median sum of education, medical, electricity, and banking availability; presence requires at least one linked village with the facility, while absence requires observed absence in every constituent village. High Female Literacy indicates an above-median share of women aged seven or older who are literate. Cutoffs use matched GPs with observed covariates before restriction to the outcome sample; missing measures remain missing.

Table~\ref{tab:het_effects} reports the interaction estimates using mean village-to-town distance, followed by sensitivity to using the minimum. Neither definition establishes a consistent relationship between development and the quota effect. The intervals permit meaningful differences.

\begin{table}[!htbp]
\centering
\caption{Development and Quota Effects: Primary Estimates and Distance Sensitivity}
\label{tab:het_effects}
\small
\setlength{\tabcolsep}{3pt}

\begin{tabular}{lllrr}
\toprule
State & Elections & Characteristic & Mean distance & Minimum distance\\
\midrule
Rajasthan & 2005--2010 & Near town & -2.3 [-7.2, 2.6] & 0.0 [-4.7, 4.6]\\
Rajasthan & 2005--2010 & High infrastructure & -0.7 [-6.5, 5.0] & -0.9 [-6.6, 4.9]\\
Rajasthan & 2005--2010 & High female literacy & 3.5 [-1.3, 8.3] & 3.2 [-1.6, 8.0]\\
Rajasthan & 2010--2015 & Near town & -0.5 [-5.4, 4.4] & -0.1 [-4.9, 4.7]\\
Rajasthan & 2010--2015 & High infrastructure & 1.6 [-3.2, 6.4] & 1.7 [-3.1, 6.5]\\
\addlinespace
Rajasthan & 2010--2015 & High female literacy & 0.5 [-4.4, 5.4] & 0.5 [-4.4, 5.4]\\
Rajasthan & 2015--2020 & Near town & 4.2 [-1.5, 9.9] & 1.6 [-3.9, 7.1]\\
Rajasthan & 2015--2020 & High infrastructure & -4.8 [-11.3, 1.8] & -4.7 [-11.2, 1.8]\\
Rajasthan & 2015--2020 & High female literacy & 5.3 [-0.4, 11.0] & 5.8 [0.0, 11.5]\\
Uttar Pradesh & 2005--2010 & Near town & 1.1 [-1.4, 3.7] & 1.2 [-1.4, 3.8]\\
\addlinespace
Uttar Pradesh & 2005--2010 & High infrastructure & 2.3 [-0.4, 4.9] & 2.3 [-0.4, 4.9]\\
Uttar Pradesh & 2005--2010 & High female literacy & 0.3 [-2.2, 2.8] & 0.3 [-2.2, 2.8]\\
Uttar Pradesh & 2010--2015 & Near town & 1.0 [-2.2, 4.2] & 1.3 [-1.9, 4.6]\\
Uttar Pradesh & 2010--2015 & High infrastructure & 0.1 [-3.0, 3.2] & 0.1 [-3.0, 3.2]\\
Uttar Pradesh & 2010--2015 & High female literacy & -1.3 [-4.5, 2.0] & -1.3 [-4.6, 2.0]\\
\addlinespace
Uttar Pradesh & 2015--2021 & Near town & -1.0 [-4.8, 2.7] & -1.8 [-5.7, 2.2]\\
Uttar Pradesh & 2015--2021 & High infrastructure & -1.7 [-5.8, 2.3] & -1.7 [-5.8, 2.3]\\
Uttar Pradesh & 2015--2021 & High female literacy & -1.1 [-4.8, 2.6] & -1.1 [-4.7, 2.6]\\
\bottomrule
\end{tabular}

\begin{minipage}{\textwidth}\footnotesize
\emph{Notes:} Interaction coefficients and 95\% block-clustered intervals are in percentage points. Each coefficient is the difference in the quota effect between the indicated group and its comparison group. All models use outcome-year block fixed effects and include the three moderators and their interactions with prior quota assignment. The columns change only the GP distance measure used to define proximity to towns. Medians use observed covariates before restriction to the outcome sample.
\end{minipage}
\end{table}
\FloatBarrier
